% ===========================================================================
%  07 -- 讨论
%
%  事实来源：第 6 章，paper_context/paper_outline.md 第 2.7 节，
%                   paper_context/system_delta_from_proposal.md
%  证据：VLM 基准测试报告，回放/安全约束证据，参与者
%            导出数据，质量门控，以及视觉事实本体演进
%  声明边界：技术可靠性与可审计性的支持力度强于教学有效性。
% ===========================================================================
\chapter{讨论}\label{ch:discussion}

第~\ref{ch:evaluation}~章报告了三条证据线，其声明强度各不相同：针对 RQ2 的视觉事实基准测试、针对 RQ1 的回放与安全约束就绪证据，以及针对 RQ3 的形成性 VR 试点。本章对这些证据线进行综合解读。核心主张并非教学系统已被证明能够改善学习效果，而是本论文展示了一种方法，通过在学习者与基础模型响应之间放置遥测数据、程序性规则、检索到的来源、VLM 导出的视觉事实以及安全约束验证，使程序性教学系统在技术上变得可审计。

第~\ref{sec:disc-evolution}~节解释了实现的工作与原提案之间的差异。
第~\ref{sec:disc-interpretation}~节对技术证据进行解读。
第~\ref{sec:disc-pilot}~节在其形成性边界内讨论试点发现。
第~\ref{sec:disc-ontology}~节从本体演进中提炼主要的设计经验。
第~\ref{sec:disc-limitations}~节和第~\ref{sec:disc-future}~节阐述了论文的局限性以及最直接的研究下一步方向。

% ===========================================================================
\section{从提案到实现系统的演变}\label{sec:disc-evolution}

原始论文提案~\citep{ThesisProposal1128} 描述了一个面向 Meta Quest~3 的 VR 优先教学系统，具备眼动追踪、手部追踪、语音输入、RAG+LLM 后端，以及一项受控的三条件人类受试者研究。最终的论文保留了沉浸式程序性教学的动机，但工作的重心发生了转移。最终实现的贡献最好理解为：一个基于证据的教学运行时系统与视觉事实提取方法，而 VR 则用于形成性试点而非受控的有效性声明。

这一转变不仅仅是工程的便利。冷启动程序使得一个前置问题变得可见：在教学系统能够提供有用的程序性帮助之前，它必须知道哪些状态可以被信任。某些所需的证据通过 DCS-BIOS 以遥测数据形式导出，某些证据仅在座舱显示屏上可见，某些步骤是手动的、不在布局中的或不适合自动确认的。因此，论文必须首先构建证据层：程序包、门控推理、证据包、检索边界、VLM 视觉事实、受 Schema 约束的模型输出、安全约束验证、叠加层白名单以及可回放日志。

VLM 工作也源于同样的压力。提案假设仿真器状态足以满足程序性跟踪的需求。然而在实践中，诸如对准和 BIT 状态等显示页内容必须通过视觉观察获取，这催生了 13 事实本体、经过人工审核的截图数据集、LoRA 适配流水线以及独立的留存基准测试。这些成为核心贡献，因为它们定义了如何将基础模型用作一个有界的感知组件，而非一个不加约束的程序状态判断器。

已完成的 VR 试点部分回应当了原始提案中人类受试者研究的宏愿。它表明，基于证据基础的运行时系统可以在头戴式设备中面向真实参与者运行，并且导出流水线可以捕获后续分析所需的数据。然而，它并未完成项目之初所提出的受控三条件研究。在 $n=9$、非随机分配且一名参与者的问卷数据未解决的情况下，该试点是形成性和描述性的。

% ===========================================================================
\section{技术证据的解读}\label{sec:disc-interpretation}

最强的技术结果是 RQ2 视觉事实基准测试。在 Qwen3.5-9B 主干上，主要 LoRA 线在两个独立的 13 事实留存集上均达到约 0.99 的事实准确率，同时相对于基座模型大幅减少了关键假阳性。补充的 Qwen3.6-27B 和 Gemma-4-31B 结果表明，相同的数据和本体配方并非绑定于单一骨干模型，虽然论文并不主张广泛的模型家族通用性。更重要的解读更窄也更有用：一个经过审核的、与本体对齐的小型数据集可以使 VLM 在座舱视觉事实提取方面的可靠性显著优于相应的基座模型（在本地留存数据上）。

这一结果的重要性在于：教学运行时系统将 VLM 输出视为证据而非权威。诸如可见页面、标签或状态文本等视觉事实是下游步骤推理的输入之一，程序引擎仍会将视觉事实与遥测数据、程序顺序、时间上下文以及程序包定义的门控结合起来。正是这种分离防止了基准测试结果沦为缺乏支撑的教学声明。VLM 结果支持教学系统的感知层，但本身并不证明教学系统的教学效果更好。

RQ1 的回放与安全约束证据支持论文的另一部分。回放案例、在线 Fixture 覆盖、程序包导出的 Schema、响应映射以及叠加层白名单表明，所实现的运行时系统能够在模型输出到达学习者之前对其进行约束。这是一种技术就绪性与可审计性的声明，意味着帮助循环可以被回放、检查和对照显式证据契约进行验证。这并不意味着每个未来的模型响应都将在教学上最优，也不意味着该运行时系统已经跨仿真器或跨程序得到了验证。

综合而言，技术证据支持了架构的核心前提：当基础模型的角色被显式证据生成与验证所约束时，它们可以在程序性教学中发挥作用。论文并非要求 LLM 作为真相的来源，而是要求 LLM 在一个运行时系统中生成响应——在该系统中，可观察状态、程序包、检索到的来源以及安全约束共同定义了响应可以安全地主张或高亮的内容。

\subsection{技术证据未能证明的内容}

技术证据未能显示学习增益、记忆保持、迁移效果、工作负荷降低或更优的训练成果。VLM 基准测试测量的是单帧视觉事实提取，回放与安全约束测试测量的是证据契约和输出约束是否按预期运行。这两者对于一个可靠的教学系统而言都是必要的，但任何一项都不足以构成人类学习效果的声明。一个学习者可能接收到准确、有据可查的帮助，却仍然无法形成持久的程序性理解；相反，一个学习者也可能从一个不那么精密但恰好契合其已有知识和策略的系统中学到东西。

证据也尚未隔离每个组件的价值。本论文未包含消融实验来对比仅遥测教学、视觉增强教学、不同检索策略或不同安全约束严格程度。因此，评估支持的是一个集成的工件声明：所实现的系统展示了一种将上述各元素在技术上可审计地组合起来的方式，而非对每个设计选择与其替代方案进行排序。

% ===========================================================================
\section{试点发现的解读}\label{sec:disc-pilot}

形成性 VR 试点提供了关于系统实际使用的描述性证据。在经过审核的数据中，所有四名有教学系统辅助的参与者完成了 33 步程序，而五名无教学系统辅助的参与者均未完成。手动编码还显示，教学系统辅助条件下无遗漏错误，而无教学系统辅助条件下的遗漏错误平均为 5.6 个。这些观察与教学系统的预期功能一致：它保持当前程序步骤可见，检索相关的程序上下文，并能指导学习者找到座舱控制。

相同的数据也表明，辅助并不能消除程序性错误。在教学系统辅助条件下，执行错误和提示辅助错误仍然存在。这是一个重要的发现，因为它指向了一个比"教学系统是否有帮助？"更精确的设计问题。剩下的问题是：灵活的模型生成式指导、确定性的安全约束修复、视觉状态的不确定性以及人类操作时机如何与严格的程序规范对齐。教学系统可以减少遗漏步骤的模式，但仍可能留下时机不当、执行不完整或过度信赖辅助的空间。

实时帮助循环日志支持了基础设施的主张。在四次有教学系统辅助的试验中，系统记录了 77 个帮助循环、23 次 VLM 调用、82 次叠加层执行、零次叠加层拒绝以及 21 次回退或修复事件。零次叠加层拒绝不应被解读为"一切顺利"。相反，它意味着在 Schema 验证和安全约束处理后，所提议的叠加层动作保持在已配置的白名单范围内。修复事件是约束层正在活跃运行的证据，而非模型输出始终直接正确的证据。

因此，该试点在论文中有两个角色。首先，它表明经过仪器化装备的教学运行时系统可以在真实 VR 环境中使用并产生可分析的日志；其次，它揭示了后续受控研究应该更仔细测量的那种人机交互数据。它不支持关于学习效果、工作负荷、记忆保持、迁移效果或任务速度的因果性主张。

% ===========================================================================
\section{本体设计经验}\label{sec:disc-ontology}

本体演进是本论文中最清晰的设计经验。初版视觉事实本体将可观察事实与程序性结论混合在一起。例如，\texttt{ins\_go} 要求 VLM 判断对准状态是否已有效达到。该目标不仅仅是视觉性的，它依赖于特定的显示文本、其他显示状态的缺失以及程序阶段。第一次基准测试中的失败模式揭示了这一边界问题：一个单帧 VLM 被要求同时执行感知和程序推理。

13 事实本体纠正了这一边界。它将复合目标分解为更低层的视觉观察，例如接地对准文本或 OK 文本是否可见，并将 FCS-MC 状态分解为视觉上可区分的阶段。下游运行时系统随后将这些观察与遥测数据和程序上下文结合起来解释。因此，设计原则是：

\begin{quote}
\textbf{视觉事实目标应该是低层级的、可局部验证的视觉命题。更高层级的程序状态应该在下游从视觉事实、遥测数据、时间上下文和程序规则中推断得出。}
\end{quote}

这条原则的意义超越了标签命名。它影响标注质量、训练信号质量、错误诊断和运行时安全性。一位人工审查者可以在不知道程序历史的情况下判断一段文本是否在截图中可见，但同一位审查者如果不知道图像之前发生了什么，则无法总是判断程序状态是否完成。当事实习契约尊重这一区别时，模型的输出变得更容易审查，也更安全地可以与确定性门控结合起来。

在完成型事实上残留的假阳性显示了该方法的当前边界。某些状态在中间阶段和最终阶段之间视觉上相似，尤其是在涉及数值读数或瞬态显示状态时。这些情况可能需要更平衡的数据、更高分辨率的裁剪、重复观察或显式引入时序信息的 VLM 输入。因此，论文并不主张 13 事实本体普遍解决了视觉接地问题，而是为一项高要求的程序提供了一种经过验证的设计纪律，以及一种用于发现视觉目标何时仍然过于接近程序性结论的具体方法。

作为面向未来领域的启发式规则，原则很简单：如果一个事实无法在不知道程序步骤的情况下从单张图像中标注，那么它对 VLM 契约来说可能层级过高，应该被分解为可见原子事实或推迟给程序引擎处理。正是在这一点上，论文从一个案例特定的实现细节转向了可复用的表示经验。

% ===========================================================================
\section{局限性}\label{sec:disc-limitations}

论文的贡献应在以下局限性范围内理解。

\subsection{评估范围}

所有数据集、回放案例和试点试验均涉及一个程序：F/A-18C 冷启动。视觉提取器在相同的机型与渲染环境下的固定复合座舱面板上训练和评估。留存集是零训练重叠的独立会话，但测量的仍为域内泛化，而非建立跨程序、跨机型或跨仿真器的通用性。

最大的训练配方在 Run-005 过采样步骤后包含 928 行双语 SFT 数据。这对于一项特定领域的适配研究而言刻意地小且具有实用价值，但它并未识别出最低数据需求、性能上限或每个配方组件的效果。论文未包含针对双语行、过采样、LoRA 秩、组合再平衡或替代微调方法的消融实验。

\subsection{系统与 VLM 范围}

VLM 作为单帧视觉事实提取器被评估。它不决定程序进度、不选择最终的教学系统动作、也不直接控制叠加层。在完成型事实上残留的关键假阳性仍然重要，因为一个错误的视觉观察如果没有被遥测数据、时序确认或安全约束检查，仍可能影响下游推理。粘性事实和确认策略缓解了这一风险，但其完整的端到端效果尚未通过实验隔离。

骨干比较涵盖了 Qwen3.5-9B、Qwen3.6-27B 和 Gemma-4-31B，在匹配的 LoRA 式适配配方下进行。这在本项目范围内是有用的模型多样性证据，但并未覆盖其他 VLM 家族、更大模型、全量微调、替代适配器或非座舱视觉领域。

运行时系统通过回放、Fixtures、桌面 DCS 和 Quest~3 试点进行了实现和演练。然而，头戴设备延迟、叠加层可读性、眼动或手部追踪精度以及专家评定的响应有用性并未作为受控的系统级结果进行测量。

\subsection{试点范围}

试点包含九名参与者：四名有教学系统辅助和五名无教学系统辅助。分配未随机化，已有 DCS 经验未平衡，样本量过小无法进行推断性统计。无教学系统辅助的试验在不同程序点停止，因此其持续时间无法与完成的有教学系统辅助试验进行比较。

问卷和工作负荷证据不完整。P02 的问卷元数据无法被可信地归因，且 NASA-TLX 整合未纳入已完成的导出流水线。因此，试点支持的是任务日志和编码观察，而非工作负荷、信心或主观可用性结论。

最后，手动审查从被动遥测门控和动作时间线推断出无教学系统辅助条件下的完成情况。由于主动 VLM 通路在该条件下不可用，视觉证据支撑的步骤确定性较低。这一局限性并不否定形成性发现，但它说明了为什么下一步研究需要更严格的仪器装备和编码。

\subsection{范围总结}

在上述局限性范围内，论文支持四项主张：一个受证据约束的程序性教学运行时系统可以以显式架构边界实现；一个小数据 VLM 适配流水线可以在本地留存集上产生可靠的视觉事实提取；视觉本体演进提供了一条关于感知与推理之间边界的具体设计经验；VR 试点表明系统可以在有参与者的条件下运行并产生可分析的日志。它并不确立受控的教育有效性。

% ===========================================================================
\section{未来工作}\label{sec:disc-future}

上述局限性指向了一个聚焦的研究议程。

\subsection{受控人类受试者评估}

下一项研究应将有依据的教学系统与基线条件进行对比，采用随机分配、平衡的参与者元数据、预定义的编码规则以及完整的问卷和工作负荷捕获。在下一轮迭代中，两条件设计相比原提案的三条件设计更为现实，因为当前的优先事项是确定已完成的运行时系统是否能在受控条件下产生可测量的完成度、错误、工作负荷和记忆保持方面的差异。

\subsection{向其他程序的本体扩展}

13 事实本体应仅在经过新的任务分析后扩展到其他程序——该分析需识别哪些步骤是遥测可见的、视觉优先的，或手动的/不在布局中的。每个新程序都需要自己的视觉事实契约、经审核的数据集、留存基准测试以及运行时绑定。这是检验本体纪律是否能迁移到冷启动之外场景的适当测试。

\subsection{时序与多帧视觉推理}

若干残留错误涉及难以从单张截图中消除歧义的状态。多帧输入、短视觉历史或重复确认门控可能有助于区分中间和最终显示状态。这需要更改捕获方式、标注、模型输入格式和基准测试指标。

\subsection{系统级帮助循环评估}

应单独评估完整的帮助循环而非各组件：触发到叠加层的延迟、高亮目标的正确性、安全约束修复的频率与类型、专家审阅下的响应有用性，以及视觉优先步骤是否优于仅遥测基线。这将更直接地将 RQ1 和 RQ2 的技术证据与 RQ3 连接起来。

\subsection{原位纠正与自适应}

当前 VLM 适配流水线是离线的。未来版本可以在视觉事实错误时收集教员或学习者的纠正，并将这些纠正输入到后续审查或增量自适应中。这一方向需要谨慎的安全防护，因为在位标签可能带有噪声，且自适应不能抹除先前可靠的行为。

\subsection{未来方向总结}

表~\ref{tab:future-work} 总结了主要的未来工作方向以及每个方向所回应的局限性。

\begin{table}[htbp]
  \centering
  \caption{未来工作方向总结。}
  \label{tab:future-work}
  \small
  \begin{tabularx}{\textwidth}{@{}>{\raggedright\arraybackslash}p{0.25\textwidth}>{\raggedright\arraybackslash}X>{\raggedright\arraybackslash}X@{}}
  \toprule
  \textbf{方向} & \textbf{当前状态} & \textbf{主要挑战} \\
  \midrule
  受控人类受试者研究 &
    形成性试点已完成；因果证据待定 &
    随机化、元数据、工作负荷与统计功效 \\

  本体扩展 &
    冷启动本体仅在当前座舱布局内进行了评估 &
    新任务分析、标注与留存设计 \\

  时序视觉推理 &
    当前 VLM 使用单张复合面板帧 &
    多帧数据、模型容量与时序指标 \\

  系统级帮助循环评估 &
    组件与回放证据已可用 &
    定义综合延迟、质量与修复指标 \\

  原位纠正与自适应 &
    离线审查与 LoRA 流水线已实现 &
    噪声标签、灾难性遗忘与安全更新策略 \\

  \bottomrule
  \end{tabularx}
\end{table}

这些方向保持了论文的核心主线不变。下一阶段不是让基础模型变得更不受约束，而是强化围绕它的证据循环：更广泛的程序包、更好的视觉事实、更丰富的时序证据、受控的人类评估，以及使每个教学系统响应在事后均可审计的日志。
